\section{System Architecture}
\label{sec:architecture}

\subsection{Overall Architecture}
Figure~\ref{fig:architecture} depicts the overall system architecture. The React frontend communicates exclusively with the FastAPI backend through a REST API; the backend is organized into the ingestion/preprocessing, model orchestration, and explanation layers introduced in Section~\ref{sec:proposed}, all three of which read from and write to a shared, file-based storage layer rather than a relational or document database, as discussed in Section~\ref{sec:proposed}.

% Overall system architecture diagram
\begin{figure}[!t]
\centering
\begin{tikzpicture}[
    node distance=0.55cm and 0.9cm,
    box/.style={rectangle, rounded corners, draw=black, thick, minimum width=2.6cm, minimum height=0.9cm, align=center, font=\scriptsize, fill=blue!8},
    dbox/.style={rectangle, rounded corners, draw=black, thick, minimum width=2.6cm, minimum height=0.9cm, align=center, font=\scriptsize, fill=green!10},
    ebox/.style={rectangle, rounded corners, draw=black, thick, minimum width=2.6cm, minimum height=0.9cm, align=center, font=\scriptsize, fill=orange!12},
    sbox/.style={rectangle, rounded corners, draw=black, thick, minimum width=5.6cm, minimum height=0.8cm, align=center, font=\scriptsize, fill=gray!8},
    arrow/.style={-{Latex[length=2mm]}, thick}
]

\node[box] (frontend) {React SPA\\Upload / Config / Results};
\node[box, below=of frontend] (api) {FastAPI REST API\\(CORS-restricted)};

\node[dbox, below left=0.5cm and -0.3cm of api] (prep) {Ingestion \&\\Preprocessing Layer};
\node[dbox, right=0.35cm of prep] (orch) {Model Orchestration\\(19-model ensemble)};
\node[ebox, right=0.35cm of orch] (xai) {Explanation Layer\\(RAG + LLM)};

\node[sbox, below=0.6cm of orch] (storage) {Disk Storage: uploads/, processed/,\\session\_data/, ChromaDB vector store};

\draw[arrow] (frontend) -- (api);
\draw[arrow] (api) -- (prep);
\draw[arrow] (api) -- (orch);
\draw[arrow] (api) -- (xai);
\draw[arrow] (prep) -- (storage);
\draw[arrow] (orch) -- (storage);
\draw[arrow] (xai) -- (storage);

\end{tikzpicture}
\caption{Overall system architecture. The React frontend communicates exclusively with the FastAPI backend over a REST API; the backend's three logical layers, ingestion/preprocessing, model orchestration, and explanation, all read from and write to a shared file-based storage layer rather than a database.}
\label{fig:architecture}
\end{figure}


\subsection{Backend Architecture}
The backend is organized as a set of FastAPI routers, one each for upload, forecasting, per-group ranking, the conversational assistant, and the business insights report, backed by a set of stateless service modules that implement preprocessing, model training and orchestration, best-model selection, per-model architectures, retrieval-augmented generation, and disk-based session persistence. This separation allows the model orchestration logic to remain entirely agnostic to which specific models accept additional arguments such as a country code or exogenous features; such per-model variation is instead handled by constructing a request-scoped override of the shared model registry at the router level.

\subsection{Machine Learning Pipeline and Prediction Flow}
Figure~\ref{fig:ml-pipeline} illustrates the training and prediction pipeline described in Section~\ref{sec:proposed}: the cleaned series is submitted to a shared thread pool that trains every applicable model concurrently, subject to the deep-learning row-count gate; every surviving model is evaluated by walk-forward cross-validation; and the composite score of Eq.~\eqref{eq:composite-score} determines the selected model and the retained leaderboard.

% Training and prediction pipeline diagram
\begin{figure}[!t]
\centering
\begin{tikzpicture}[
    node distance=0.4cm and 0.5cm,
    stepbox/.style={rectangle, rounded corners, draw=black, thick, minimum width=3.0cm, minimum height=0.75cm, align=center, font=\scriptsize, fill=blue!8},
    modelbox/.style={rectangle, draw=black, thick, minimum width=1.55cm, minimum height=0.6cm, align=center, font=\tiny, fill=orange!10},
    arrow/.style={-{Latex[length=2mm]}, thick}
]

\node[stepbox] (upload) {Uploaded file\\(CSV / Excel / JSON)};
\node[stepbox, below=of upload] (clean) {Clean, gap-filled,\\outlier-capped series};
\node[stepbox, below=of clean] (fanout) {Submit to thread pool\\(one task per model)};

\node[modelbox, below left=0.5cm and 2.1cm of fanout] (stat) {Statistical\\(7 models)};
\node[modelbox, right=0.25cm of stat] (ml) {Machine\\Learning (4)};
\node[modelbox, right=0.25cm of ml] (dl) {Deep Learning (7)\\[-2pt]{\tiny gated by row count}};
\node[modelbox, right=0.25cm of dl] (croston) {Intermittent\\(CrostonSBA)};

\node[stepbox, below=1.0cm of ml] (cv) {Walk-forward CV per model\\(adaptive fold count)};
\node[stepbox, below=of cv] (score) {Composite score $S$\\(Eq.~\ref{eq:composite-score}, Tukey-fence normalized)};
\node[stepbox, below=of score] (select) {Select minimum-$S$ model;\\retain full leaderboard};
\node[stepbox, below=of select] (respond) {Forecast + interval +\\leaderboard returned to UI};

\draw[arrow] (upload) -- (clean);
\draw[arrow] (clean) -- (fanout);
\draw[arrow] (fanout) -- (stat);
\draw[arrow] (fanout) -- (ml);
\draw[arrow] (fanout) -- (dl);
\draw[arrow] (fanout) -- (croston);
\draw[arrow] (stat) -- (cv);
\draw[arrow] (ml) -- (cv);
\draw[arrow] (dl) -- (cv);
\draw[arrow] (croston) -- (cv);
\draw[arrow] (cv) -- (score);
\draw[arrow] (score) -- (select);
\draw[arrow] (select) -- (respond);

\end{tikzpicture}
\caption{Training and prediction pipeline. All nineteen registered models are submitted concurrently to a shared thread pool; deep learning models are withheld when the series falls below the frequency-specific minimum row count. Every surviving model is evaluated via walk-forward cross-validation and ranked by the composite score of Eq.~\eqref{eq:composite-score}.}
\label{fig:ml-pipeline}
\end{figure}


\subsection{Data Flow}
Figure~\ref{fig:data-flow} details the automated data-cleaning flow applied to every uploaded dataset, exactly as executed for the evaluation dataset described in Section~\ref{sec:dataset}.

% Data preprocessing flow diagram
\begin{figure}[!t]
\centering
\begin{tikzpicture}[
    node distance=0.45cm,
    stepbox/.style={rectangle, rounded corners, draw=black, thick, minimum width=6.6cm, minimum height=0.65cm, align=center, font=\scriptsize, fill=green!8},
    arrow/.style={-{Latex[length=2mm]}, thick}
]

\node[stepbox] (s1) {1. Column detection (date / target / group / exogenous)};
\node[stepbox, below=of s1] (s2) {2. Date parsing cascade \& target numeric coercion};
\node[stepbox, below=of s2] (s3) {3. Index set, sorted; duplicate timestamps summed};
\node[stepbox, below=of s3] (s4) {4. Frequency inferred; calendar gaps reindexed};
\node[stepbox, below=of s4] (s5) {5. Missing values: interpolate $\to$ ffill $\to$ bfill};
\node[stepbox, below=of s5] (s6) {6. IQR outlier detection and capping};
\node[stepbox, below=of s6] (s7) {7. Exogenous features aligned (price: mean/ffill; promo: max/zero)};
\node[stepbox, below=of s7, fill=blue!10] (s8) {Clean target series + aligned exogenous frame};

\foreach \a/\b in {s1/s2, s2/s3, s3/s4, s4/s5, s5/s6, s6/s7, s7/s8}
    \draw[arrow] (\a) -- (\b);

\end{tikzpicture}
\caption{Automated data preprocessing pipeline applied to every uploaded dataset prior to model training, exactly as executed for the evaluation dataset in Section~\ref{sec:dataset} (Table~\ref{tab:preprocessing-steps}).}
\label{fig:data-flow}
\end{figure}


\subsection{Deployment Architecture}
Figure~\ref{fig:deployment} shows the deployment topology. The frontend is served as a static single-page application build, and the backend runs as an independent ASGI process; the two communicate over HTTP with cross-origin resource sharing restricted to an explicit allow-list, extensible for a production frontend origin via an environment variable. The only external network dependency in the entire architecture is the large-language-model API used exclusively by the explanation layer; the ingestion, preprocessing, and forecasting functionality of the system operates fully without external network access.

% Deployment architecture diagram
\begin{figure}[!t]
\centering
\begin{tikzpicture}[
    node distance=0.5cm and 0.6cm,
    box/.style={rectangle, rounded corners, draw=black, thick, minimum width=2.7cm, minimum height=0.8cm, align=center, font=\scriptsize, fill=blue!8},
    ext/.style={rectangle, rounded corners, draw=black, thick, minimum width=2.7cm, minimum height=0.8cm, align=center, font=\scriptsize, fill=red!8, dashed},
    arrow/.style={-{Latex[length=2mm]}, thick}
]

\node[box] (browser) {Client browser};
\node[box, below=of browser] (vite) {Vite-served React\\static bundle};
\node[box, below=of vite] (uvicorn) {Uvicorn ASGI process\\(FastAPI application)};
\node[ext, right=1.1cm of uvicorn] (llm) {Gemini API\\(chat / embeddings)};
\node[box, below=of uvicorn] (disk) {Local disk:\\uploads, processed,\\session\_data, ChromaDB};

\draw[arrow] (browser) -- (vite);
\draw[arrow] (browser) -- (uvicorn);
\draw[arrow] (uvicorn) -- (disk);
\draw[arrow] (uvicorn) -- (llm);

\end{tikzpicture}
\caption{Deployment architecture. The frontend and backend are independently deployable processes; CORS origins and the frontend's backend base URL are environment-configured. The only external network dependency is the large-language-model API used by the explanation layer; core forecasting functions without it.}
\label{fig:deployment}
\end{figure}

